%---
% slug: marey
% title: Marey
% category: product
% eyebrow: SCIENTIFIC VISUALIZATION
% summary: A plotting library for F# where the compiler checks your units, and the axes label themselves. Scales, SI prefixes, ticks, and axis titles are derived from unit witnesses, and the same figure renders to SVG, PNG, PDF, or LaTeX source.
% author: Jonathan Schäfer
% draft: false
% date: 2026-08-21
% rank: 4
%---
\documentclass[11pt]{article}
\usepackage{publication-web}
\begin{document}

\publicationtitle{Marey, a plotting library for F\# where the compiler checks your units}
\publicationauthors{Jonathan Schäfer}
\publicationaffiliations{Axym Labs, Darmstadt, Germany}

\begin{publicationhead}
\makepublicationtitle
\end{publicationhead}

\publicationfigure[0.9\linewidth]
  {figures/marey-free-fall.png}{/work/marey-free-fall.svg}
  {Free fall: displacement against elapsed time}
  {Free fall: displacement against elapsed time. Nothing configures an axis.}

\publicationabstract{A plotting library for F\# where the compiler checks your units, and the axes label themselves. Marey is named for \'Etienne-Jules Marey (1830--1904), who invented the graphical method in experimental science: record a physical quantity as a curve against time, and read the physics off the curve.}

\publicationactions{%
  \publicationaction{https://github.com/Jonathan-Schaefer-git/marey}{View code}
  \publicationaction{https://www.nuget.org/packages/Marey.Plot}{Open package}
  \publicationaction{https://github.com/Jonathan-Schaefer-git/marey/blob/main/docs/guide.md}{Read the guide}%
}

\begin{verbatim}
#r "nuget: Marey.Plot"
#r "nuget: Marey.Backend.Svg"

open Marey

let t = Array.init 24 (fun i -> float i * 0.4<s>)
let x = t |> Array.map (fun t -> 0.5 * 9.81<m/s^2> * t * t)

Chart.line U.second U.metre t x
|> Chart.withTitle "Free fall"
|> Chart.toSvg SvgBackend.render "fall.svg" (Chart.size 800.0 500.0)
\end{verbatim}

Nothing above configures an axis. The scale, the domain, the SI prefix, the tick values, the tick labels and both axis titles are derived from the two unit witnesses. \texttt{x} is never annotated as a length --- \texttt{0.5 * 9.81<m/s\^{}2> * t * t} has type \texttt{float<m>}, which the compiler derives --- and plotting seconds against seconds on these axes is a compile error.

\section{Install}

\begin{verbatim}
dotnet add package Marey.Plot
dotnet add package Marey.Backend.Svg
\end{verbatim}

For PNG and PDF, add \texttt{Marey.Backend.Skia}, the one package with third-party dependencies.

\section{Units, end to end}

Data spanning $10^{-9}$ to $5 \cdot 10^{-9}$ metres produces an axis reading \texttt{nm} with ticks 1, 2, 3, 4, 5. Kilopascals are labelled as kilopascals:

\publicationfigure[0.9\linewidth]
  {figures/marey-wave-packet.png}{/work/marey-wave-packet.svg}
  {A Gaussian wave packet in millimetres and kilopascals}
  {A Gaussian wave packet in millimetres and kilopascals.}

The unit registry is keyed on the seven-element SI dimension vector rather than on strings, so anything that composes to a newton prints \texttt{N}. Units with no SI name are covered as long as the composition is stated: volts divided by metres render as volts per metre, and as megavolts per metre when prefixed. Nothing is inferred: a bare dimension is never factorised, because the same vector is both \texttt{J/K} and \texttt{N$\cdot$m/K}.

Affine units carry their offset: degrees Celsius converts by shift as well as factor, an electronvolt axis reads \texttt{keV} rather than \texttt{kJ}, and a mass rescaled by $10^{-6}$ prints milligrams rather than micro-kilograms.

\section{Figures that belong in a LaTeX document}

Swapping the theme turns the same figure into \texttt{tikzpicture} source. It is typeset by the same engine as the rest of the document, so \textbf{the figure inherits the document's fonts}, and the axis titles are set by \texttt{siunitx} from the unit witnesses:

\begin{verbatim}
figure
|> Chart.withTheme Theme.tikz
|> Chart.toFile TikzBackend.render "fall.tex" (Chart.size 320.0 200.0)
\end{verbatim}

\section{The usual scientific figure}

A model curve, measured points, error bars, and a legend naming only the layers that were named:

\publicationfigure[0.9\linewidth]
  {figures/marey-oscillator.png}{/work/marey-oscillator.svg}
  {A damped oscillator: model against measurement with error bars}
  {A damped oscillator: model against measurement with error bars.}

\section{When the data is a table}

When the data is a table and the figure has more than two things in it, \texttt{plot} says what each column means in one place:

\begin{verbatim}
plot trials {
    x U.second (fun r -> r.Time)
    y U.metre  (fun r -> r.Displacement)
    color      (fun r -> r.Sample)
    layer GeomPoint
    fit
}
\end{verbatim}

\publicationfigure[0.9\linewidth]
  {figures/marey-warming-rate.png}{/work/marey-warming-rate.svg}
  {Global temperature anomaly by era, with a least-squares trend}
  {Global temperature anomaly by era, with a least-squares trend.}

Each dimensioned channel takes its witness beside the accessor, and the accessor is checked against it: naming \texttt{r.Displacement} under \texttt{U.second} does not compile. Categorical channels take no witness, because a category carries no dimension.

\texttt{fit} is one of the statistics, and its slope is a typed number as well as a line on the page. A fit of metres against seconds has a slope of type \texttt{float<m/s>}:

\begin{verbatim}
let fitted = Stat.regression U.second U.metre t x
let speed: float<m/s> = fitted.Slope
\end{verbatim}

\texttt{facetWrap} draws one panel per level of a key over shared axes, so a difference between panels is a difference in the data.

\section{Real data, and where the units story stops}

Units are transformative for physical measurements and irrelevant for a mole fraction in ppm, which is dimensionless. A dimensionless axis is rendered with no unit brackets, values with no unit to carry a prefix are not rescaled, and the label is left to the caller:

\publicationfigure[0.9\linewidth]
  {figures/marey-co2-record.png}{/work/marey-co2-record.svg}
  {Atmospheric CO2 at Mauna Loa, annual means with uncertainties}
  {Atmospheric CO\textsubscript{2} at Mauna Loa, annual means with uncertainties.}

The repository's \texttt{examples/} directory has twelve worked scripts, three of them on real committed datasets.

\section{A figure that moves}

An animation is a \textbf{pure function of time}, so seeking is evaluation: \texttt{anim.At 2.7<s>} is a direct call, frames come out in any order, and a render is \texttt{Array.Parallel.map} with no scheduler. The storyboard block writes the sequence:

\begin{verbatim}
open Marey.Storyboard

let intro =
    scene {
        let! curve = draw measured |> over 1.2<s>
        do! wait 0.5<s>
        do! parallel' [ curve |> morphTo fitted |> over 2.0<s>
                        curve |> zoomToY (0.0, 0.045) |> over 2.0<s> ]
        do! highlight curve.[peak]
    }

let anim = animate 32 (Chart.size 320.0 220.0) intro
\end{verbatim}

\publicationfigure[\linewidth]
  {figures/marey-storyboard.png}{/work/marey-storyboard.svg}
  {Six frames of one storyboard: a ringdown drawn, morphed onto its envelope, zoomed, and one mark picked out}
  {Six frames of one storyboard: a ringdown drawn, morphed onto its envelope, zoomed, and one mark picked out.}

\texttt{let!} binds a \textbf{name} for something on stage, which is what lets \texttt{curve} be referred to three more times. The whole shot is laid out \textbf{once}, against the union of every state sampled, so the panel holds still while the data moves rather than crawling sideways as the tick labels gain a digit.

Rendering is a map over the frame times, on as many threads as there are. The sequential and concurrent passes produce the same bytes in the same order, and a gate checks this frame for frame on both backends.

\section{What is in the box}

\begin{itemize}
\tightlist
\item \texttt{Marey.Core}: dimension vectors, unit witnesses, scales, and the ten-case scene IR. Zero dependencies.
\item \texttt{Marey.Units}: SI, 22 derived and 18 non-SI unit witnesses, prefix selection, and Unicode, ASCII, LaTeX and \texttt{siunitx} rendering.
\item \texttt{Marey.Plot}: \texttt{Chart} and the \texttt{plot} block, twelve geoms, statistics, facets, six themes, colour-blind-safe palettes, the \texttt{scene} storyboard block, and the parallel frame writer.
\item \texttt{Marey.Backend.Svg}, \texttt{Marey.Backend.Tikz}, \texttt{Marey.Backend.Skia}: SVG, \texttt{tikzpicture} source, and PNG and vector PDF.
\item \texttt{Marey.Anim}: animation as a seekable, composable function of time, with twenty-eight easings and arc-length path morphing.
\item \texttt{Marey.Tools.Video}, \texttt{Marey.Tools.Tex}, \texttt{Marey.Toolchain}: video encoding through \texttt{ffmpeg} and TeX typesetting through \texttt{latex} and \texttt{dvisvgm}, both degrading gracefully when the program is absent.
\item \texttt{Marey.Html}, \texttt{Marey.Verso}: notebook display, so a bare \texttt{Figure} shows as a figure and a \texttt{Film} plays as a looping flipbook.
\end{itemize}

\section{Tests}

The suite combines FsCheck property suites for the dimension group laws, all seven scale kinds, ticks, interpolation, colour and the scene IR; golden baselines in SVG and \texttt{tikzpicture} source; a real \texttt{pdflatex} compile gate; rasterised ink gates over the compiled TeX, the rendered PNG and the rendered PDF; and executable checks over every example in the guide and in \texttt{examples/}.

Several of those gates shell out, and each passes when the program it needs is absent. So each run prints a census of all nine external dependencies and what became of them, and \texttt{MAREY\_REQUIRE\_TOOLCHAIN=1} turns an absent one into a failure. CI sets it.

\section{Status}

Pre-1.0, and versioned accordingly. The API is still free to move; breaking changes are listed in the changelog with the reason and the replacement. MIT licensed.

\end{document}
